%% First draft (ICML 2025 style). Derived from ../main.tex (progress report), which stays unchanged.
\documentclass{article}
\usepackage{microtype}
\usepackage{graphicx}
\usepackage{booktabs}
\usepackage{multirow}
\usepackage{xcolor}
\usepackage{hyperref}
\newcommand{\theHalgorithm}{\arabic{algorithm}}
\usepackage[accepted]{icml2025}   % non-anonymous, as the course allows
\usepackage{amsmath,amssymb}
\usepackage[capitalize,noabbrev]{cleveref}

% Placeholder markers: `grep -n "todo\|tbd" paper.tex` must be empty before submitting.
\newcommand{\todo}[1]{\textcolor{red}{[TODO: #1]}}
\newcommand{\tbd}{\textcolor{red}{??}}
\newcommand{\fig}[2]{\IfFileExists{#1}{\includegraphics[width=#2]{#1}}%
  {\fbox{\parbox[c][3cm][c]{#2}{\centering\small\todo{\detokenize{#1}}}}}}

\icmltitlerunning{Where vs.\ What: Landmarks and Pixels for Low-Data Sign Recognition}

\begin{document}
\twocolumn[
\icmltitle{Where vs.\ What: Hand Landmarks and Pixels\\for Low-Data Isolated Sign Recognition}
\begin{icmlauthorlist}
\icmlauthor{Narelle Ho}{ucb}
\icmlauthor{Cheng Yao Lee}{ucb}
\end{icmlauthorlist}
\icmlaffiliation{ucb}{CS289A, University of California, Berkeley}
\icmlcorrespondingauthor{Cheng Yao Lee}{\todo{email}}
\icmlkeywords{Sign language recognition, WLASL, pose estimation, video classification}
\vskip 0.3in
]
\printAffiliationsAndNotice{}

%=====================================================================
\begin{abstract}
%=====================================================================
Isolated sign recognition maps a short video of one sign to a word. Two paradigms compete: classifiers on
hand landmarks from an off-the-shelf pose estimator, and video networks that learn from pixels. We
compare them on exactly the same 1{,}013 WLASL100 clips (7.5 training clips per sign), scoring every model
on accuracy, training time, and inference latency over three seeds. On landmarks, a 167k-parameter
transformer reaches 72.3\% validation and 66.7\% test top-1, and tree, MLP, and recurrent ladders show that
temporal aggregation, not architecture, drives accuracy. On pixels, a from-scratch CNN stays at chance and a
frozen Kinetics-pretrained R(2+1)D reaches only 20\%; fine-tuned, it reaches 37--43\% validation. To
separate knowing \emph{where} the hands are from representing \emph{what} shape they make, we feed the
backbone landmark-guided hand crops: localization changes accuracy by at most 5 points, within noise,
leaving a gap of about 30 points that only the skeletal representation closes. Once landmarks are
extracted, the transformer is 500$\times$ faster at inference with 188$\times$ fewer parameters, but
extraction itself costs about 2\,s per clip on a CPU. We release code and discuss the limits of a
half-sized benchmark with overlapping signers.
\end{abstract}

%=====================================================================
\section{Introduction}
%=====================================================================
More than 70 million deaf people use a sign language \citep{wfd}, yet most speech and captioning tools
assume a spoken utterance. Isolated sign recognition, naming the single sign performed in a short clip,
is a narrow, measurable version of this gap and a building block for continuous recognition
\citep{koller2020survey}.

Two families of methods dominate. \emph{Pose-based} methods run a pretrained keypoint detector such as
MediaPipe \citep{lugaresi2019mediapipe} and classify the resulting trajectory
\citep{li2020wlasl,bohacek2022spoter,selvaraj2022openhands}. \emph{Appearance-based} methods learn from
RGB frames with video backbones pretrained on action recognition \citep{carreira2017quo,tran2018closer}.
Pose models are cheap and robust to background; pixel models see everything, including cues the detector
misses. Published comparisons usually differ in data, preprocessing, and budget at the same time, so it
is hard to tell \emph{why} one paradigm wins.

We ask a narrower question in the low-data regime, where it matters most: \emph{is the landmark
advantage due to localizing the hands or to the skeletal abstraction itself?} Our contributions:
\begin{itemize}
\item A controlled benchmark on the 1{,}013 WLASL100 clips still obtainable, with three-seed means,
      training time, latency, and parameters for five model families on identical splits (\cref{sec:setup}).
\item Ablation ladders for each pose model showing that temporal pooling or attention and augmentation,
      not architecture, account for most of the accuracy (\cref{sec:pose}).
\item A 2$\times$2 pixel study (frozen vs.\ fine-tuned backbone $\times$ full frame vs.\ landmark-guided hand
      crop) that isolates the value of localization (\cref{sec:pixels}).
\item An analysis of feature separability, signer leakage, and errors that qualifies what the
      absolute numbers mean (\cref{sec:analysis}).
\end{itemize}

%=====================================================================
\section{Related Work}
%=====================================================================
\textbf{Word-level sign datasets.} WLASL \citep{li2020wlasl} collects 21{,}083 videos of 2{,}000 ASL
signs from the web; its 100-sign subset WLASL100 is a common benchmark. The original paper compares
appearance (I3D) and pose (Pose-GRU, Pose-TGCN) baselines, reporting 65.9\%, 46.5\%, and 55.4\% top-1 on
WLASL100. The Google ISLR competition \citep{kaggle2023islr} released landmark-only data at a much larger
scale; we use WLASL because it provides raw video, which makes a pixel comparison possible.

\textbf{Pose-based recognition.} Keypoint trajectories from OpenPose or MediaPipe
\citep{zhang2020mediapipehands} are classified with GRUs \citep{li2020wlasl}, graph convolutions
\citep{li2020wlasl,tunga2021gcnbert}, and transformers \citep{bohacek2022spoter}. OpenHands
\citep{selvaraj2022openhands} standardizes pose pipelines across languages, and SignBERT
\citep{hu2021signbert} pretrains a hand-model-aware encoder. These works scale model capacity; we instead
hold the representation fixed and ablate aggregation and augmentation across four model families,
including a tree ensemble \citep{chen2016xgboost}, which remains competitive on tabular data
\citep{grinsztajn2022tree}.

\textbf{Appearance-based recognition.} Inflated 3D CNNs \citep{carreira2017quo} and factorized
R(2+1)D convolutions \citep{tran2018closer} pretrained on Kinetics \citep{kay2017kinetics} are standard
backbones. \citet{li2020transferring} transfer knowledge from subtitled news signs; SAM-SLR
\citep{jiang2021skeleton} fuses RGB, depth, and skeleton streams and crops around the hands. Unlike these
systems, we do not fuse modalities: we use the landmark detector only to \emph{place} a crop, so pixels and
landmarks share localization but differ in representation.

%=====================================================================
\section{Setup}\label{sec:setup}
%=====================================================================
\subsection{Data}
We take the first 100 glosses of \texttt{WLASL\_v0.3.json} with the official split, which reproduces the
published counts (2{,}038 videos; 1{,}442/338/258). Many source links are dead; the preprocessed Kaggle
copy contains \textbf{1{,}013} of these videos. Every model in this paper is trained and scored on that
same subset and split (\cref{tab:data}). With 7.5 training clips per sign, overfitting is the main
difficulty. The 100-clip test set covers 72 signs and carries roughly $\pm$5 points of sampling error,
so we select every configuration on validation accuracy and report test accuracy only for the selected
ones. 31 of the 34 test signers also appear in training, as in the official split.

\begin{table}[t]
\centering\footnotesize\setlength{\tabcolsep}{3pt}
\begin{tabular}{lcc}
\toprule
 & Full WLASL100 & Available subset \\
\midrule
Videos (tr / val / test) & 1{,}442 / 338 / 258 & 748 / 165 / 100 \\
Train clips / sign (min/mean/max) & 12 / 14.4 / 30 & 4 / 7.5 / 12 \\
Signs present in test & 100 & 72 \\
Signers (tr / val / test) & 91 / 69 / 56 & 60 / 41 / 34 \\
Test signers also in train & 52 of 56 & 31 of 34 \\
Median length (10--90th pct.) & -- & 65 (33--95) frames \\
Frames with L / R / both hands & -- & 32\% / 60\% / 27\% \\
\bottomrule
\end{tabular}
\caption{Dataset statistics. Hand detection by MediaPipe Hands.}
\label{tab:data}
\end{table}

\subsection{Landmark representation}
MediaPipe Hands \citep{zhang2020mediapipehands} gives 21 3D keypoints per hand per frame; we map its
mirrored handedness labels to the signer's own left and right. Per frame and hand we center the points on
the wrist and divide by palm length (wrist to middle-finger base), add a missing-hand flag, and
optionally the wrist position in the frame, giving 132 values per frame (128 without location). Clips are
resampled to $T{=}32$ frames. Sequence models read the $32\times132$ sequence; the tree and MLP read it
flattened, or pooled over time (mean and standard deviation, 264 values, or means of four chunks).
Features are standardized with training statistics. Training-time augmentation applies one random
temporal crop (up to 10\% per end), rotation ($\pm15^\circ$), scale ($\pm10\%$), shift ($\pm5\%$), and
point jitter per clip.

\subsection{Models}
\textbf{Gradient-boosted trees.} XGBoost \citep{chen2016xgboost} with depth-3 trees, learning rate 0.08,
row/column subsampling 0.8/0.3, $\lambda{=}2$, up to 300 rounds with early stopping on validation, and
five augmented copies of the training set.
\textbf{MLP and linear probe.} Feed-forward networks on the flattened or pooled vector.
\textbf{GRU / LSTM} \citep{cho2014gru,hochreiter1997lstm}. 128 hidden units, ablated over direction,
readout (last state vs.\ temporal mean), and depth.
\textbf{Transformer} \citep{vaswani2017attention}. A linear projection to $d{=}128$, LayerNorm, learned
positional embeddings, 4 heads, feed-forward width 256, dropout 0.3, and a last-token or mean readout.
\textbf{From-scratch CNN.} Three conv--BN--ReLU blocks on $8\times64\times64$ frames, averaged over time.
\textbf{Pretrained video backbone.} R(2+1)D-18 \citep{tran2018closer} pretrained on Kinetics-400
\citep{kay2017kinetics} (31.4M parameters), reading $T\times112\times112$ clips. We either freeze it and
train a linear head on its 512-d pooled features, or fine-tune it end to end with the same augmentation
applied to pixels (learning rates $10^{-4}$ backbone, $10^{-3}$ head, batch 8, fp16, at most 30 epochs
with patience 6).
The \emph{full-frame} input letterboxes each frame to a square; the \emph{hand-crop} input cuts one square
per clip around the 2nd--98th percentile of all detected hand keypoints, enlarged by 25\%
(\cref{fig:crops}). The crop uses only the \emph{location} of the landmarks, never their values.

\textbf{Training and evaluation.} Neural models use AdamW (weight decay $10^{-2}$), label smoothing 0.1,
and early stopping on validation. Every configuration is run with seeds 0, 1, 2; we report the mean
validation and test top-1, test top-5, and the seed standard deviation, the mean training time, and
batched inference latency per clip. All models run on one RTX 3060 GPU (XGBoost on CPU). Latency
counts model inference only; it excludes video decoding for both paradigms and MediaPipe extraction for
the landmark models, which we measure separately at 30\,ms per frame on CPU (about 2\,s for a median
65-frame clip).

\begin{figure}[t]
\centering
\fig{figs/crops.png}{\columnwidth}
\caption{Full-frame (left) vs.\ landmark-guided hand-crop (right) inputs to the video backbone, four of
16 frames per clip. In 16:9 videos the hands occupy a few pixels of a $112\times112$ full frame.}
\label{fig:crops}
\end{figure}

%=====================================================================
\section{Results}
%=====================================================================
\subsection{Main comparison}
\cref{tab:main} reports the configuration of each family selected on validation, and
\cref{fig:frontier} places them on an accuracy--latency plane. The landmark transformer is the most
accurate model on validation (72.3\%) and on test (66.7\%), 29.5 validation points above the best pixel
model, a fine-tuned R(2+1)D on hand crops. Given cached landmarks, it is also the cheapest: 188$\times$ fewer
parameters, 57$\times$ shorter training, and 500$\times$ lower inference latency than the backbone. The
caveat is extraction: MediaPipe Hands costs about 2\,s of CPU time per clip, two orders of magnitude more
than the backbone's 12\,ms on a GPU, so the landmark pipeline is cheaper only when extraction is amortized
or run on-device.

\begin{table*}[t]
\centering\footnotesize
\begin{tabular}{llrcccrr}
\toprule
Family & Input & Params & Val top-1 & Test top-1 & Test top-5 & Train (s) & Latency (ms/clip) \\
\midrule
XGBoost (4-chunk pooled)       & landmarks & $\le$30k trees & 47.5 & 39.3 $\pm$ 1.2 & 73.7 & 86 & 0.13$^\dagger$ \\
MLP (mean/std pooled)$^\ast$   & landmarks & --   & 63.6 & 60.0 & 84.3 & $<$60 & -- \\
GRU (bidir., mean)$^\ast$      & landmarks & --   & 66.3 & 60.7 & 87.7 & $<$60 & -- \\
LSTM (bidir., mean)$^\ast$     & landmarks & --   & 63.4 & 63.3 & 86.3 & $<$60 & -- \\
Transformer (1 layer, mean)    & landmarks & 167k & \textbf{72.3} & 66.7 $\pm$ 1.9 & 86.3 & 29 & 0.02$^\dagger$ \\
\midrule
CNN from scratch (16 frames)   & pixels    & 107k & 3.4 & 1.3 $\pm$ 0.5 & 3.3 & 197 & 0.34$^\dagger$ \\
R(2+1)D frozen, hand crop      & pixels + loc. & 31.4M$^\ddagger$ & 19.6 & 12.3 $\pm$ 1.9 & 22.3 & 13 & 12.0 \\
R(2+1)D fine-tuned, full frame & pixels    & 31.4M & 37.4 & \textbf{38.7} $\pm$ 3.1 & 65.0 & 2034 & 11.8 \\
R(2+1)D fine-tuned, hand crop  & pixels + loc. & 31.4M & \textbf{42.8} & 34.3 $\pm$ 3.4 & \textbf{65.7} & 1683 & 11.5 \\
\midrule
\multicolumn{8}{l}{\emph{Reference, full WLASL100 (2{,}038 clips) \citep{li2020wlasl}:} Pose-GRU 46.5, Pose-TGCN 55.4, I3D 65.9 test top-1.} \\
\bottomrule
\end{tabular}
\caption{Best configuration per family, selected on validation; mean over 3 seeds ($\pm$ seed std of test
top-1). $^\ast$From the MLP and RNN ladders (\cref{tab:pooled,tab:seq}). $^\dagger$Measured on the
default configuration. $^\ddagger$51k trainable. Latency is batched model inference only; MediaPipe
extraction adds about 30\,ms per frame on CPU to every landmark model.}
\label{tab:main}
\end{table*}

%---------------------------------------------------------------------
\subsection{Pose models: aggregation beats architecture}\label{sec:pose}
%---------------------------------------------------------------------
\begin{table}[t]
\centering\footnotesize\setlength{\tabcolsep}{2.5pt}
\begin{tabular}{lcccccc}
\toprule
 & \multicolumn{2}{c}{Linear} & \multicolumn{2}{c}{MLP} & \multicolumn{2}{c}{XGBoost} \\
\cmidrule(lr){2-3}\cmidrule(lr){4-5}\cmidrule(lr){6-7}
Rung & Val & Test & Val & Test & Val & Test \\
\midrule
1. Shape, 32 frames flat & 27.7 & 16.0 & 36.6 & 28.3 & 38.8 & 31.3 \\
2. + wrist location      & 31.1 & 18.7 & 39.8 & 30.3 & 40.4 & 37.3 \\
3. + augmentation        & 43.4 & 36.0 & 52.5 & 44.0 & 46.9 & 45.7 \\
4a. + pool (mean, std)   & 52.3 & 52.7 & \textbf{63.6} & \textbf{60.0} & 39.6 & 41.3 \\
4b. + pool (4 chunks)    & \textbf{56.2} & 45.0 & 62.2 & 53.0 & \textbf{47.5} & 39.3 \\
5. 4a + velocity         & 48.9 & 46.0 & 57.0 & 53.3 & 37.6 & 40.3 \\
\bottomrule
\end{tabular}
\caption{Ladder for models that read a fixed-length vector; top-1 (\%), mean over 3 seeds (seed std
$\le$3.4). Bold: selected on validation.}
\label{tab:pooled}
\end{table}

\begin{table}[t]
\centering\footnotesize\setlength{\tabcolsep}{2.5pt}
\begin{tabular}{lcccccc}
\toprule
 & \multicolumn{2}{c}{GRU} & \multicolumn{2}{c}{LSTM} & \multicolumn{2}{c}{Transformer} \\
\cmidrule(lr){2-3}\cmidrule(lr){4-5}\cmidrule(lr){6-7}
Rung & Val & Test & Val & Test & Val & Test \\
\midrule
0. Base, no augmentation   & 41.8 & 43.0 & 41.6 & 47.3 & 70.7 & 68.0 \\
1. Base (1 layer, last)    & 51.3 & 53.3 & 50.5 & 50.0 & 71.9 & 70.0 \\
2. + bidirectional         & 57.0 & 53.3 & 57.0 & 51.0 & -- & -- \\
3. + mean over time        & \textbf{66.3} & \textbf{60.7} & \textbf{63.4} & \textbf{63.3} & \textbf{72.3} & \textbf{66.7} \\
4. + second layer          & 66.3 & 65.7 & 60.0 & 56.0 & 70.1 & 67.0 \\
5. 4 without location      & -- & -- & -- & -- & 67.3 & 62.7 \\
\bottomrule
\end{tabular}
\caption{Ladder for sequence models on the $32\times132$ input; top-1 (\%), mean over 3 seeds (seed std
$\le$5.0). Bold: selected on validation (ties go to the smaller model).}
\label{tab:seq}
\end{table}

\cref{tab:pooled,tab:seq} add one design choice per rung.
\textbf{Temporal aggregation is the largest single gain} for models that cannot align time themselves:
mean/std pooling lifts the MLP from 44.0\% to 60.0\% test, and a temporal-mean readout lifts the GRU from
53.3\% to 60.7\%. Models that see the whole sequence through attention are far less sensitive: the
transformer is within noise of its best configuration even at the base rung without augmentation
(70.7\% val).
\textbf{Augmentation matters most for flat inputs} (+14--17 points for linear and MLP, +8 for XGBoost) and
least for the transformer (+1.2 val).
\textbf{Trees do not benefit from pooling}: XGBoost peaks on the augmented flat vector (45.7\% test),
and pooling loses 4--6 test points, likely because axis-aligned splits on 100 classes with 7.5 clips per
class cannot exploit the dense, correlated pooled features that a linear layer combines easily.
\textbf{Once a model aggregates over time, architecture matters little}: the selected MLP, GRU, and LSTM
are within 3.3 test points, inside the $\pm$5-point test noise; the transformer is best on validation by
6 points over the next model.
The selected neural pose models all exceed the original Pose-GRU and Pose-TGCN baselines trained on
twice the data \citep{li2020wlasl}, though the protocols differ (\cref{sec:limits}).

%---------------------------------------------------------------------
\subsection{Pixel models: where vs.\ what}\label{sec:pixels}
%---------------------------------------------------------------------
\begin{table}[t]
\centering\footnotesize\setlength{\tabcolsep}{3pt}
\begin{tabular}{lccccr}
\toprule
Rung & Val & Test & Top-5 & Train (s) \\
\midrule
CNN from scratch, $8\times64^2$       & 3.2 & 1.0 & 3.7 & 37 \\
\quad + augmentation                  & 3.0 & 1.7 & 5.0 & 80 \\
\quad + 16 frames                     & 3.4 & 1.3 & 3.3 & 197 \\
\midrule
R(2+1)D frozen, full, 8 frames        & 18.4 & 11.7 & 20.3 & 8 \\
R(2+1)D frozen, full, 16 frames       & 18.0 & 10.7 & 19.7 & 14 \\
R(2+1)D frozen, hand crop, 16         & 19.6 & 12.3 & 22.3 & 13 \\
R(2+1)D fine-tuned, full, 16          & 37.4 & \textbf{38.7} & 65.0 & 2034 \\
R(2+1)D fine-tuned, hand crop, 16     & \textbf{42.8} & 34.3 & \textbf{65.7} & 1683 \\
\bottomrule
\end{tabular}
\caption{Pixel models; top-1 (\%), mean over 3 seeds (seed std $\le$3.4). Chance is 1\%. Frozen
training time includes feature extraction.}
\label{tab:pixels}
\end{table}

The from-scratch CNN stays at chance in every configuration (\cref{tab:pixels}): resizing a full frame to
$64\times64$ blurs the hands into the torso and background, and 748 clips are too few to learn what to
look at. A Kinetics-pretrained backbone changes this, but only when fine-tuned.
\textbf{Frozen action features do not encode signs}: a linear head on frozen R(2+1)D features reaches
18--20\% validation, and doubling the frames from 8 to 16 does not help (18.4 vs.\ 18.0), so temporal
resolution is not the bottleneck. \textbf{Fine-tuning doubles accuracy} to 37--43\% validation and 65\%
test top-5, at a cost of about 30 GPU-minutes per seed, against 30 seconds for the transformer.

\textbf{Where vs.\ what.} If the landmark advantage came from localization, handing the backbone the
hand region should close most of the gap. It does not. With a frozen backbone the crop adds 1.6 validation
points; fine-tuned, it adds 5.4 validation points but \emph{loses} 4.4 test points, and pooled over all
265 held-out clips the two inputs differ by under 2 points (39.6\% vs.\ 37.9\%). Against a 29.5-point
validation gap to the transformer, localization explains at most about a sixth of the difference. The rest
comes from representing the hand as a normalized skeleton: invariant to skin, clothing, lighting, and
camera, and already disentangled into the shape features that our probes find carry most of the signal
(\cref{sec:analysis}). Two caveats bound this conclusion. Two of six fine-tuning runs were still improving
at the 30-epoch cap, so the pixel numbers are a lower bound; and our crop is one fixed box per clip that
still contains the face and torso (\cref{fig:crops}), not per-frame hand patches.

\begin{figure}[t]
\centering
\fig{figs/acc_vs_latency.png}{\columnwidth}
\caption{Test top-1 vs.\ inference latency (log scale) for the selected configuration of each family;
error bars are seed std.}
\label{fig:frontier}
\end{figure}

%=====================================================================
\section{Analysis}\label{sec:analysis}
%=====================================================================
\textbf{What the landmarks encode.} We probe the hand features with nearest-neighbor and logistic
regression classifiers on val+test (265 clips; \cref{tab:probe}). Time-averaged features beat the flat
sequence (41.1\% vs.\ 27.9\%), because signers start and move at different times, which is exactly the
gain temporal aggregation delivers in \cref{sec:pose}. Hand shape alone reaches 26.0\%; location alone
12.1\%, and adds only 1.9 points to shape. The features also encode the signer: the nearest training
clip shares the signer (26.0\%) more often than the sign (22.3\%).

\begin{table}[t]
\centering\footnotesize\setlength{\tabcolsep}{3pt}
\begin{tabular}{lclc}
\toprule
Probe & Acc. & Probe & Acc. \\
\midrule
1-NN / 5-NN (cos.) & 22.3 / 15.8 & Linear, shape only & 26.0 \\
Linear, flattened & 27.9 & Linear, location only & 12.1 \\
Linear, time-avg. & \textbf{41.1} & NN same signer (chance) & 26.0 (5.1) \\
\bottomrule
\end{tabular}
\caption{Feature separability on val+test (265 clips; \%, chance 1\%).}
\label{tab:probe}
\end{table}

\textbf{Errors.} Taking each model's majority prediction over three seeds, the transformer misclassifies
32 of the 100 test clips and the hand-crop backbone 69. The landmark errors are nearly a subset of the
pixel errors: 25 of the transformer's 32 failures are also backbone failures, so the clips that defeat
landmarks are hard for pixels too, while the backbone additionally fails on clips that landmarks solve.
The transformer's most frequent confusion, \textsc{before}$\to$\textsc{enjoy} (6 of 9 possible errors,
\cref{fig:confusions}) is consistent with two signs whose hand shapes are similar and which differ mainly
in movement relative to the body, information that wrist-centered shape features do not carry. The backbone's confusions are spread
thinly (at most 3 per pair) across visually unrelated signs, consistent with a model that has not yet
learned a stable notion of hand shape from 748 clips. \todo{Narelle: check 2--3 example clips for the top
pairs (occlusion, out-of-frame hands, failed detection).}

\begin{figure*}[t]
\centering
\fig{figs/confusions.png}{0.8\textwidth}
\caption{Most frequent test confusions (true $\to$ predicted), summed over three seeds, for the selected
landmark transformer and the selected pixel model (fine-tuned R(2+1)D on hand crops).}
\label{fig:confusions}
\end{figure*}

%=====================================================================
\section{Conclusion}\label{sec:limits}
%=====================================================================
On the 1{,}013 obtainable WLASL100 clips, a 167k-parameter transformer on MediaPipe hand landmarks is the
strongest model we trained, and temporal aggregation and augmentation, rather than architecture, explain
most of the spread between pose models. A from-scratch CNN cannot learn from 7.5 clips per sign; a
fine-tuned pretrained video backbone can, but stays about 30 validation points behind. Giving the backbone
the hand location through a landmark-guided crop recovers at most a few of those points, so in this
low-data regime the landmark advantage comes mainly from \emph{what} the skeleton represents, not from
knowing \emph{where} to look. Once landmarks are cached, that representation is also about 500$\times$
cheaper at inference; extraction, at about 2\,s per clip on CPU, is the dominant cost of the pipeline.

\textbf{Limitations.} (i) \emph{Half the data:} only 1{,}013 of 2{,}038 clips are available, so absolute
numbers are not comparable to published WLASL100 results, and the test set covers 72 of 100 signs.
(ii) \emph{Small test set:} 100 clips give about $\pm$5 points of sampling error; we select on
validation and report seed variation, but differences under $\sim$5 points between families are not
significant. (iii) \emph{Signer overlap:} 31 of 34 test signers appear in training, and the features
encode the signer, so accuracy on unseen signers is likely lower. (iv) \emph{Budgets differ:} the backbone
had a 30-epoch cap that two runs reached while still improving, one backbone, and a single crop policy;
MLP and RNN rows come from separate runs without latency measurements.

\textbf{Future work.} Signer-disjoint splits; body and face landmarks (MediaPipe Holistic) for location
relative to the body; fusing the hand-crop backbone with landmarks; and recovering the missing videos
from their original sources.

\textbf{Contributions.} Cheng Yao Lee: data loader, XGBoost, transformer, video backbones, comparison, paper.
Narelle Ho: keypoint extraction, feature analysis, features and augmentation, MLP, GRU/LSTM.
Code: \url{https://github.com/chengyaolee/asl-benchmark}.

\bibliography{refs}
\bibliographystyle{icml2025}

\end{document}
